\documentclass{article}
\usepackage{graphicx} % Required for inserting images
\usepackage{amsmath, amssymb}

\title{Trading Horizon as Endogenous Decision}
\author{Thiago, Phi}
\date{May 2026}

\begin{document}

\maketitle

\section{Theory}

\subsection{Trading Horizon as an Endogenous Decision}

Consider an investor trading a collection of assets indexed by
$i \in \{1, \ldots, N\}$. At any decision time $t$, the investor may
express a position over one of several admissible trading horizons. Let
$
    \mathcal{H} = \{h_1, h_2, \ldots, h_K\}
$
denote a finite, discrete set of horizons. A horizon is not just a
sampling convention, it defines the interval over which a trading
decision is intended to earn a return and over which its risk and
implementation cost are realized.

For each horizon $h \in \mathcal{H}$, let
$q_{i,t}^{(h)} \in [-1,1]$ denote the signed fraction of the capital
assigned to horizon $h$ that is exposed to asset $i$ at time $t$.
Positive values represent long positions and negative values represent
short positions. Thus the choice is real-valued: the investor may assign
partial long or short exposure to an asset rather than making only an
all-or-nothing decision.

The investor has finite available capital, normalized to one. Let
$a_{h,t} \in [0,1]$ denote the fraction of total capital allocated to
trading decisions at horizon $h$. Capital allocations must satisfy

\[
    \sum_{h \in \mathcal{H}} a_{h,t} \leq 1,
    \qquad
    \sum_{i=1}^{N} \left|q_{i,t}^{(h)}\right| \leq 1
    \quad \text{for each } h \in \mathcal{H}.
\]

The absolute-value restriction bounds gross exposure within each
horizon allocation, allowing offsetting long and short positions without
implicitly permitting unlimited leverage. Any unallocated share may be
understood as capital held outside the risky positions. The construction
of the candidate exposures is left unrestricted at this stage. The
resulting signed fraction of total capital exposed to asset $i$ is

\[
    Q_{i,t} = \sum_{h \in \mathcal{H}} a_{h,t} q_{i,t}^{(h)}.
\]

Let $r_{i,t,t+h}$ be the return of asset $i$ over the interval associated
with horizon $h$, and let $C_{t,t+h}^{(h)}$ represent the costs incurred
by implementing the corresponding horizon-specific allocation. The net
payoff attributable to horizon $h$ is

\[
    \Pi_{t,t+h}^{(h)}
    =
    a_{h,t}
    \sum_{i=1}^{N} q_{i,t}^{(h)} r_{i,t,t+h}
    -
    C_{t,t+h}^{(h)}.
\]

Therefore the investor's decision is not only which assets to trade,
but also which trading horizons should receive capital. Distinct
horizons may offer different expected payoffs, risks, costs, and
dependence with one another. An allocation that ignores the horizon
dimension implicitly assumes that the same trading window is always
appropriate, or that the relative value of alternative windows is
constant.

The basic hypothesis of this framework is that trading horizon is an
endogenous portfolio choice. Conditional on the information available
at time $t$, the allocation vector
$
    a_t = (a_{h,t})_{h \in \mathcal{H}}
$
is chosen jointly with the real-valued asset exposure choices, subject
to the investor's finite capital and risk constraints. The objective is to allocate exposure to
the set of horizon-specific opportunities that provides the most
attractive net risk-adjusted payoff. 

\subsection{Decision Times and Overlapping Horizon Positions}

Let $\delta>0$ denote the shortest interval at which the investor may
revise the portfolio. Decision times lie on the common grid
$
    \mathcal{T}
    =
    \{t_0+n\delta:n=0,1,\ldots,M\}.
$
The admissible horizons are aligned with this grid:
$
    \mathcal{H} \subseteq \{m\delta:m\in\mathbb{N}\}.
$
Consequently, for every $t\in\mathcal{T}$ and every
$h\in\mathcal{H}$, both the initiation time $t$ and the maturity time
$t+h$ lie on the investor's decision grid.

At each decision time $t$, the investor may initiate a new position for
every admissible horizon, including a new longer-horizon position even
when earlier positions at that same horizon remain active. Define the
newly initiated signed exposure to asset $i$ for the position vintage
$(t,h)$ as
$
    x_{i,t}^{(h)} = a_{h,t}q_{i,t}^{(h)}.
$
A vintage initiated at time $t$ for horizon $h$ remains active until
$t+h$. Thus, positions initiated at different decision times can have
overlapping payoff intervals and distinct maturity times. At any
decision time $s\in\mathcal{T}$, total active exposure to asset $i$ is
\[
    Q_{i,s}
    =
    \sum_{h \in \mathcal{H}}
    \sum_{\substack{t \in \mathcal{T} \\ t \leq s < t+h}}
    x_{i,t}^{(h)}.
\]

Two notions of capacity are relevant when horizon-specific positions
overlap. First, let $G^{\mathrm{exec}}>0$ denote the admissible gross
exposure of the position actually executed in the market after active
vintages are netted within each asset. Feasible executed exposure
satisfies
\[
    \sum_{i=1}^{N} \left|Q_{i,s}\right| \leq G^{\mathrm{exec}},
    \qquad s \in \mathcal{T}.
\]
Second, let $B^{\mathrm{commit}}>0$ denote the permitted aggregate
commitment to all active horizon-specific decisions before netting.
This is not a broker-level capital constraint; it is an overlap-control
restriction that limits simultaneous, potentially conflicting,
horizon-specific exposure. Active decisions must satisfy
\[
    \sum_{i=1}^{N}
    \sum_{h \in \mathcal{H}}
    \sum_{\substack{t \in \mathcal{T} \\ t \leq s < t+h}}
    \left|x_{i,t}^{(h)}\right|
    \leq B^{\mathrm{commit}},
    \qquad s \in \mathcal{T}.
\]
With normalized capital and no leverage, $G^{\mathrm{exec}}$ may be set
to one; $B^{\mathrm{commit}}$ is chosen as a separate modeling limit.
The two restrictions serve different purposes: the first limits the
broker-level portfolio after offsetting positions are combined, while
the second prevents opposing position vintages from being added without
limit merely because they net to a small executed exposure.
Thus overlapping horizon decisions do not provide new capital. A
position associated with a longer horizon continues to consume active
decision commitment capacity when newer decisions are made.

\subsection{The Dynamic Allocation Problem}

Let $W_t$ denote investor wealth immediately before the decision at time
$t$, and let
$
    x_t =
    \left(x_{i,t}^{(h)}\right)_
    {\substack{i=1,\ldots,N \\ h \in \mathcal{H}}}
$
collect the new horizon-specific exposures selected at that time. The
investor chooses $x_t$ conditional on the information available at
$t$ and on all position vintages that remain active from earlier
decisions. In particular, $x_t$ may include a newly initiated
longer-horizon position even if positions with the same horizon were
initiated at earlier decision times and have not yet matured.
In general form, the decision problem is
\[
    x_t^*
    \in
    \underset{x_t \in \mathcal{A}_t}{\operatorname{arg\,max}}
    \left\{
        \mathbb{E}_t\!\left[\Delta W(x_t)\right]
        -
        \gamma \mathcal{R}_t(x_t)
        -
        \mathcal{C}_t(x_t;Q_t^{\mathrm{existing}})
    \right\},
\]
where $\mathbb{E}_t[\Delta W(x_t)]$ is the conditional expected payoff
from the decision, $\mathcal{R}_t(x_t)$ is its risk contribution,
$\mathcal{C}_t(x_t;Q_t^{\mathrm{existing}})$ is the implementation and
adjustment cost given the existing portfolio, and $\gamma \geq 0$
controls risk aversion.

The feasible set enforces both capacity constraints after combining the
new decision with active exposures. Let
$\mathcal{V}_t^{\mathrm{existing}}$ be the set of position vintages
initiated prior to $t$ that have not matured. Then
\[
    \mathcal{A}_t
    =
    \left\{
        x_t :
        \sum_{i=1}^{N}
        \left|
            Q_{i,t}^{\mathrm{existing}}
            +
            \sum_{h \in \mathcal{H}}x_{i,t}^{(h)}
        \right|
        \leq G^{\mathrm{exec}},
        \quad
        \sum_{i=1}^{N}
        \left(
            \sum_{(\tau,h)\in\mathcal{V}_t^{\mathrm{existing}}}
            \left|x_{i,\tau}^{(h)}\right|
            +
            \sum_{h\in\mathcal{H}}\left|x_{i,t}^{(h)}\right|
        \right)
        \leq B^{\mathrm{commit}}
    \right\}.
\]
The investor may consequently revise the allocation across trading
horizons at every decision time, but only within the portfolio's
executed-exposure and active-decision commitment limits and after
accounting for active positions, risk, and adjustment cost.

\end{document}
